\documentclass[11pt]{article}

\usepackage[margin=1in]{geometry}
\usepackage{amsmath,amssymb,amsthm}
\usepackage{booktabs}
\usepackage{enumitem}
\usepackage{hyperref}
\usepackage{microtype}
\usepackage{algorithm}
\usepackage{algpseudocode}

\hypersetup{hypertexnames=false,colorlinks=true,linkcolor=blue,citecolor=blue,urlcolor=blue}

\newtheorem{theorem}{Theorem}
\newtheorem{lemma}{Lemma}
\newtheorem{remark}{Remark}

\title{\textbf{Disagreement-Gated Recertification and Boundary Audits for Anonymous DHT Routers}\\
\large Keeping certified anonymity--latency menus valid under churn}
\author{}
\date{March 23, 2026 (archive v683)}

\begin{document}
\maketitle

\begin{abstract}
Anonymous DHT defenses are rarely static: operators change retry rules, cover-traffic rates, and thresholded ``attack-mode'' triggers as churn and adversary pressure shift.
Re-certifying every tweak via fresh, fully labeled privacy audits is often prohibitive.

We give a practical recertification protocol that preserves the many-testing-safe guarantees of Certified~A while spending audit effort only where the new policy differs from the old.
The core device is a \emph{zero-label disagreement gate}: if the new policy's disagreement rate with a previously certified policy is below the old certificate margin, then the update is immediately safe.
If not, we show how to label only the disagreement region and how to concentrate boundary audits where tail/stop-time constraints are tight.
A small DHT-shaped toy table (kept in-\TeX) illustrates audit savings. This note is intentionally the certified-routing \emph{change-control companion} beneath Certified~A and Certified~B rather than a competing runtime-monitoring or release paper. Its stop point is the recertification packet: same-claim check $\rightarrow$ zero-label gate $\rightarrow$ disagreement-band certificate $\rightarrow$ boundary-audit schedule, with runtime conformance receipts deferred to Operational~B and evaluator/public packaging deferred to Eval~1 and ABOM/OINL~+~Release~A.
\end{abstract}


\paragraph{Scope.}
We address \emph{recertification under policy drift}: how to keep certified anonymity/latency menus valid when operators tune retries, thresholds, and cover parameters. The paper is intentionally positioned as the certified-routing change-control companion beneath Certified~A's theorem/accountant root and Certified~B's optional multiplicity-compression manifest, not as a second runtime-monitoring or release-lineage entrypoint.

\paragraph{Imports.}
We import certified-menu objects and the baseline many-testing guarantees from Certified~A, and routing-signature quantizers from Certified~B.\cite{series:menus,series:sigcomp}
We also import stop-time / tail-normal-form background from the published wiki links referenced by Certified~A.
Exposure-normal-form vocabulary (what counts as the \emph{same} transcript interface) is imported from the trace-interfaces note.\cite{series:traceifaces}
Recertification outputs are contract objects (Synthesis~0) and should be anchored by ABOM/OINL manifests when published.\cite{series:contractobjects,series:abomoinl}
Receipt line-item claim fields (including \texttt{tw\_id} and optional \texttt{state\_decl\_id}) are standardized in Synthesis~12; replay/state drift rules are centralized in Synthesis~18.\cite{series:receiptitems,series:planstateequiv}
Runtime conformance receipts for a pinned deployment plan are imported from Operational~B, while evaluator notarization and release-facing comparison are imported from Eval~1 and Release~A.\cite{series:operationalB,series:eval1,series:releaseA}

\paragraph{Exports.}
Recertification packet: claim identity, zero-label gate, disagreement audit rule, and boundary-audit schedule. Publicly, it answers only three questions: whether the update stays under the old claim object, whether a small disagreement-band packet closes the gap, or whether the update is large enough to require a new claim object.

\paragraph{Boundary.}
This paper stops once a fixed certified interface has produced a recertification decision object for a proposed update: either ``new claim object required'' or ``same claim, with gate verdict and any incremental disagreement-band / boundary-audit evidence.'' The maintained worked cut keeps the needed source objects visible as 
\nolinkurl{example\_change\_control.json}, 
\nolinkurl{example\_compare\_report.json}, 
\nolinkurl{example\_drift\_cases.json}, and 
\nolinkurl{example\_release\_action\_matrix.json}; it does \emph{not} resolve a separate 
\nolinkurl{bundle://worked-example/recertification} in 
\nolinkurl{example\_support\_bundle\_map.json}, does not add an 
\nolinkurl{eval-recertification-v1} plan row, and does not emit a separate recertification packet in 
\nolinkurl{example\_verifier\_report.json}. This paper therefore remains the first public owner of the change-control packet while the worked example supplies only a source-object boundary for the guard / classification / case-ladder drill. Continuous runtime conformance receipts belong to Operational~B, evaluator notarization belongs to Eval~1, and public lineage / release packaging belong to ABOM/OINL plus Release~A.\cite{series:workedexample,series:operationalB,series:eval1,series:releaseA} In particular, this note does not re-own the live deployment-conformance evidence needed to show that the declared policy actually ran, nor the signed publication receipt that exposes the resulting claim to the world.

\paragraph{Earliest sufficient stop.}
If a verifier can inspect the ENF/TW/state identity verdict, the disagreement upper bound, the zero-label gate verdict, and any incremental disagreement-band / boundary-audit packet, then every extra claim owned by this paper is already fixed. In this cut the maintained worked example exposes that verdict only through the change-control, compare-report, drift-case, and release-action source objects, not through a materialized recertification support-bundle route. Runtime plan-pinning / epoch receipts remain with Operational~B, while notarization and signed public lineage remain with Eval~1 and ABOM/OINL~+~Release~A.\cite{series:workedexample,series:operationalB,series:eval1,series:releaseA}

\paragraph{Artifact policy.}
Source-only; full scripts and raw sizing curves are available on request.\cite{series:artifactpolicy}

\paragraph{Later-terse citation posture.}
Later terse papers should cite this note only when the live issue is the \emph{change-control packet itself}: whether the update is the same claim, what the disagreement upper bound is, whether the zero-label gate passed, and what incremental disagreement-band / boundary-audit packet closes any remaining gap.
If the live issue is instead the menu theorem/accountant root, multiplicity compression, runtime plan conformance, evaluator notarization, or signed public lineage, stop at Certified~A, Certified~B, Operational~B, Eval~1, or Release~A respectively rather than widening this paper into a second theorem root or release note.
If the live issue is instead the concrete worked recertification boundary, stop at this note plus 
\nolinkurl{example\_change\_control.json}, 
\nolinkurl{example\_compare\_report.json}, 
\nolinkurl{example\_drift\_cases.json}, and 
\nolinkurl{example\_release\_action\_matrix.json}; do not cite a non-existent 
\nolinkurl{bundle://worked-example/recertification} route in this archive cut. Widen beyond that boundary only when runtime conformance, evaluator notarization, or signed release lineage has become the live question.\cite{series:menus,series:sigcomp,series:workedexample,series:operationalB,series:eval1,series:releaseA}

\section{Problem: continuous tuning in Anonymous DHTs}
Anonymous DHT deployments rarely hold a routing policy fixed.
Operators adjust retry caps, cover ratios, and thresholded ``attack-mode'' branches as churn, load, and adversary pressure shift.
Certified~A defines \emph{certified menus} (and many-testing-safe certificates) for a baseline policy $\pi_0$; this paper addresses the \emph{recertification} problem:
given a proposed update $\pi_1$, when can we accept it \emph{without} paying for a full new labeled privacy audit?

Our guiding principle is to spend work only where the update changes behavior.
We contribute:
(i) a \emph{zero-label disagreement gate} that accepts an update when its disagreement rate with $\pi_0$ is below the old certificate margin;
(ii) a targeted procedure that labels only the disagreement region when the gate fails; and
(iii) \emph{boundary audits} that concentrate effort near tight tail/stop-time constraints.
Throughout, we treat ``policy'' broadly: it may be a concrete router, a choice from a certified menu, or a dial-sheet instantiation of a budgeted knob set.
For DHT background, menu objects, and the baseline certification story, see Certified~A~\cite{series:menus}.

\section{Setup and notation}
Fix an SLA witness $W_\pi$ encoded as a mean inequality $\mathbb{E}[W_\pi]\le 0$ as in Certified~A.
Assume the witness is bounded and, for the update pair $(\pi_0,\pi_1)$, satisfies $|W_{\pi_1}-W_{\pi_0}|\le 1$ (true for Bernoulli-offset SLAs like privacy, coverage, and tail indicators, and for bounded losses in $[0,1]$).
Let $\pi_0$ be a currently deployed policy with a certificate margin:
\begin{equation}
\mathbb{E}[W_{\pi_0}] \le -\gamma
\qquad (\gamma>0).
\label{eq:margin}
\end{equation}
We consider a proposed update $\pi_1$ and want to decide whether (and how) to recertify it.

\section{First question: is this the same claim?}
Before invoking any disagreement gate, a deployment must decide whether the update preserves the \emph{certified interface}.
Certified~A certificates and the gate below are statements about a declared witness surface, repetition unit, and public conditioning.
If those move, the update is not merely ``a new policy on the same claim''; it may be a \emph{different claim object}.\cite{series:traceifaces,series:contractobjects}

\paragraph{Receipt-level identity check.}
In this archive, ``same claim'' is intended to be mechanical: the public receipt tuple includes \texttt{exposure\_nf\_id} (the witness interface), \texttt{tw\_id} (the repetition window/threat declaration), and optionally \texttt{state\_decl\_id} (the state surface assumed).
If any of these ids changes, the safe default is to mint a new claim object (and then decide whether it is a conservative coarsening of a prior interface).\cite{series:receiptitems,series:planstateequiv}

\begin{lemma}[Interface-preserving drift is a prerequisite for the gate]
\label{lem:sameclaimfirst}
Consider a baseline certificate for policy $\pi_0$ at exposure normal form $(\mathcal Y,g,\mathsf{unit},\mathsf{cond})$.
A recertification shortcut based on disagreement between $g_{\pi_0}$ and $g_{\pi_1}$ is sound only if one of the following holds:
\begin{enumerate}[leftmargin=*]
\item the update keeps the same exposure normal form up to metadata/version labels; or
\item the updated certificate explicitly coarsens the old interface, i.e.\ the new observation is $g' = h\circ g$ under the same repetition unit and public conditioning.
\end{enumerate}
If the update changes the observation map in any other way, or changes the repetition unit or public conditioning, then the correct action is to mint a new claim object and certify that new interface directly.
\end{lemma}

\begin{proof}[Proof sketch]
The disagreement argument compares two policies on the \emph{same experiment}: the same witness extractor, composed over the same repetition unit and under the same public conditioning.
If the interface changes only by an explicit coarsening, the certified endpoint can only improve by data processing, so reuse is conservative.
Otherwise the witness channel itself changes, and the original accountant no longer certifies the updated claim.
\end{proof}

\paragraph{Operational rule.}
Run a ``same exposure normal form?'' check \emph{before} the zero-label gate.
Only after the interface is fixed should a deployment spend unlabeled evidence on disagreement estimation.

\paragraph{Aside: replay-plan drift.}
If the privacy claim is unchanged but the replay workflow changes (\texttt{plan\_spec\_id} drift), the publication must still republish a new receipt/ABOM pinning the updated replay promise.
Whether that drift triggers new privacy evidence is classified by the P0/P1/P2 taxonomy in Synthesis~18.\cite{series:planstateequiv}

\section{Zero-label disagreement gate}
Let $g_\pi(X)$ denote the \emph{routing decision} induced by policy $\pi$ in context $X$ (which template is used; which threshold branch; \ldots).
For randomized policies, treat the per-episode PRG seed $U$ as part of the context; see Certified~B~\cite{series:sigcomp}.
In implementation, it is often simplest to set $g_\pi(X):=Q(\pi(X))$ for the same quantizer $Q$ used in routing-signature compression (Certified~B), so that disagreement is measured on \emph{realized} decisions.
In particular, $g_\pi$ should include any action components that can affect an SLA witness (padding schedule, retry count, dummy rate), not just a coarse ``mode bit.''
Define the \emph{disagreement rate}
\[
q := \mathbb{P}\bigl(g_{\pi_0}(X)\neq g_{\pi_1}(X)\bigr),
\]
estimable from unlabeled logs (contexts only).
With a finite unlabeled log of size $n_u$, we estimate $q$ by $\widehat{q}$ (the empirical disagreement rate) and form a one-sided upper confidence bound $\overline{q}$.
As in Certified~A, compute a one-sided upper confidence bound $\overline{q}$ on $q$ from unlabeled contexts (e.g., the exact Clopper--Pearson upper bound) \cite{clopper1934}.
This is the same structural object that governs disagreement-based active learning (the ``disagreement region''), though our goal is certification rather than learning \cite{hanneke2014}.

\paragraph{Closed form for OR-threshold routers (uniform context model).}
If $X=(a,c)$ is uniform on $[0,1]^2$ and $g_{\pi}(a,c)=\mathbf{1}\{a>\theta_a \ \text{or}\ c>\theta_c\}$, then $g_{\pi}=0$ only on the rectangle $[0,\theta_a]\times[0,\theta_c]$.
The disagreement rate between $\pi_0=(\theta_a,\theta_c)$ and $\pi_1=(\theta_a',\theta_c')$ is the symmetric-difference area of these rectangles:
\begin{equation}
q = \theta_a\theta_c + \theta_a'\theta_c' - 2\min(\theta_a,\theta_a')\min(\theta_c,\theta_c').
\label{eq:disagree-area}
\end{equation}
In practice $q$ is estimated from unlabeled logs under the deployment context distribution (not necessarily uniform), but the rectangle formula is useful for sanity checks and for designing threshold step sizes that stay within a desired disagreement budget.

We impose a mild condition that holds for most thresholded anonymous routing designs:
\begin{quote}
Whenever two policies take the same routing decision on a context, their SLA witness is identical on that episode.
Formally, $W_{\pi_0}(X,Y)=W_{\pi_1}(X,Y)$ whenever $g_{\pi_0}(X)=g_{\pi_1}(X)$.
\end{quote}

\begin{lemma}[Zero-label safety check from disagreement]
\label{lem:disagree}
Assume $W_{\pi_0}(X,Y)=W_{\pi_1}(X,Y)$ whenever $g_{\pi_0}(X)=g_{\pi_1}(X)$ and $|W_{\pi_1}-W_{\pi_0}|\le 1$.
If $\pi_0$ satisfies \eqref{eq:margin}, then
\[
\mathbb{E}[W_{\pi_1}] \le -\gamma + q.
\]
In particular, if $q\le \gamma$, then $\mathbb{E}[W_{\pi_1}]\le 0$ and the update is safe \emph{without any new labels}.
\end{lemma}

\begin{remark}[Multiple SLAs: gate by the tightest margin]
If Certified~A certified $K$ SLA witnesses with margins $\gamma_1,\dots,\gamma_K$ for $\pi_0$, then the same disagreement rate $q$ upper-bounds the change for each bounded witness.
Thus a sufficient zero-label condition to preserve \emph{all} SLAs is $q\le \min_k \gamma_k$ (apply Theorem~\ref{thm:gate-est} with a union bound over the $K$ gates).
If only a subset of SLAs is near-tight, it is often effective to gate on those and immediately trigger targeted audits for the rest.
\end{remark}

\paragraph{Interpretation.}
If the new policy disagrees with the old only rarely (under the context distribution), then the old certificate ``pays'' for the change.

\paragraph{Relation to safe policy improvement.}
The gate is a DHT-native analogue of \emph{safe policy improvement}: we only accept an update when we can certify it does not violate a previously certified baseline constraint.
In reinforcement-learning and bandit deployments, closely related ideas accept policy updates based on high-confidence off-policy evaluation and baseline guarantees \cite{thomas2015hcope,laroche2019spibb}.
Our setting differs in that the witness is a per-episode, bounded SLA functional (privacy/coverage/tail/cost) and we exploit the discrete routing structure to localize audits to a disagreement region.


\begin{theorem}[Gate with estimated disagreement (finite-sample)]
\label{thm:gate-est}
Assume $\pi_0$ has been certified so that \eqref{eq:margin} holds with probability at least $1-\delta_0$ under the deployment distribution.
Let $\overline q$ be any upper confidence bound on the disagreement rate $q$ computed from an \emph{independent} unlabeled context log, such that $\mathbb{P}(q\le \overline q)\ge 1-\delta_1$ (e.g., Clopper--Pearson inversion).
If the update is accepted only when $\overline q\le \gamma$, then with probability at least $1-(\delta_0+\delta_1)$, the accepted update satisfies $\mathbb{E}[W_{\pi_1}]\le 0$.
\end{theorem}

\section{When the gate fails: label only the disagreement region}
If $q>\gamma/2$, we should not blindly deploy, but we still do not need to label everything.
Let $\mathcal{R}=\{X: g_{\pi_0}(X)\neq g_{\pi_1}(X)\}$ be the disagreement region.
Only episodes with $X\in\mathcal{R}$ can change the SLA witness.

Writing $\Delta W := W_{\pi_1}-W_{\pi_0}$, note that $\Delta W=0$ off $\mathcal{R}$, so
\[
  \mathbb{E}[W_{\pi_1}] = \mathbb{E}[W_{\pi_0}] + q\,\mathbb{E}[\Delta W \mid X\in\mathcal{R}].
\]
Thus, when the gate fails ($q>\gamma$), it suffices to label only $X\in\mathcal{R}$ and obtain a one-sided upper bound on $\mathbb{E}[\Delta W \mid X\in\mathcal{R}]$, yielding a tight incremental certificate.

\begin{theorem}[Incremental certificate from disagreement-band labels]
\label{thm:incremental}
Assume $\pi_0$ satisfies \eqref{eq:margin} with probability at least $1-\delta_0$.
Let $\overline q$ be an upper confidence bound on $q$ from an independent unlabeled log such that $\mathbb{P}(q\le \overline q)\ge 1-\delta_1$.
Let $\overline{\Delta}$ be an upper confidence bound on the conditional mean $\mathbb{E}[\Delta W\mid X\in\mathcal{R}]$ computed from (possibly selectively sampled) labeled data on $\mathcal{R}$ with known propensities (e.g., via Horvitz--Thompson / IPS weighting), such that $\mathbb{P}(\mathbb{E}[\Delta W\mid X\in\mathcal{R}]\le \overline{\Delta})\ge 1-\delta_2$.
Then with probability at least $1-(\delta_0+\delta_1+\delta_2)$,
\[
\mathbb{E}[W_{\pi_1}] \le -\gamma + \overline q\,\overline{\Delta}.
\]
In particular, if $\overline q\,\overline{\Delta}\le \gamma$ then the update is certified-safe.
\end{theorem}


\paragraph{Importance correction for selective labeling.}
If we only label a subset of episodes (e.g., only those with $X\in\mathcal{R}$), the labeled sample is biased.
A clean deployment pattern is to assign a known labeling propensity $\rho(X)\in(0,1]$ concentrated on $\mathcal{R}$ (e.g., $\rho(X)=\rho_0$ on $\mathcal{R}$ and a small floor $\rho_{\min}>0$ elsewhere if unconditional estimates are needed), and apply inverse-probability weighting.
Let $L_i\sim\mathrm{Bernoulli}(\rho(X_i))$ be the label-indicator and let $\widetilde{W}_i$ be the witness computed when the label is obtained.
Then the Horvitz--Thompson / IPS estimator
\[
\widehat{\mathbb E}[W_{\pi_1}]\;:=\;\frac{1}{n}\sum_{i=1}^n \frac{L_i\,\widetilde{W}_i}{\rho(X_i)}
\]
is unbiased under mild conditions \cite{horvitz1952}.
If the deployment monitors or re-audits continuously, the statistical meaning should survive optional stopping and repeated publication.
For a receipt-friendly (union-bound) way to publish a single monitoring confidence budget and ``spend'' it over repeated audit events, see Synthesis~26.\cite{series:anytimespend}

A minimal recertification protocol:
\begin{enumerate}
  \item estimate $q$ from unlabeled contexts;
  \item if $q\le \gamma$, accept (Lemma~\ref{lem:disagree});
  \item else, selectively label episodes in $\mathcal{R}$ and certify $\mathbb{E}[W_{\pi_1}\mid X\in\mathcal{R}]$ with an importance correction.
\end{enumerate}

\begin{algorithm}[t]
\caption{Disagreement-gated recertification (single SLA)}
\label{alg:gate}
\begin{algorithmic}[1]
\Require Certified policy $\pi_0$ with margin $\gamma$; proposed update $\pi_1$; unlabeled contexts; overall failure $\delta$.
\State Estimate $\widehat{q}$ from contexts and form an upper confidence bound $\overline{q}=\mathrm{CP\_UCB}(\widehat{q};n_u,\delta)$ using the one-sided Clopper--Pearson (exact binomial) bound, where $n_u$ is the unlabeled context sample size.
\If{$\overline{q}\le \gamma$}
  \State \Return \textbf{Accept} update (zero labels).
\Else
  \State Collect/label episodes with a known propensity $\rho(X)$ concentrated on $\widehat{\mathcal{R}}$ (disagreement band).
  \State Certify $\mathbb{E}[W_{\pi_1}]\le 0$ using IPS/HT weighting on labeled data, with Bonferroni/Holm correction as needed.
  \State \Return \textbf{Accept} only if certified.
\EndIf
\end{algorithmic}
\end{algorithm}

\noindent\textbf{Implementation note.} $\mathrm{CP\_UCB}(\widehat{q};n_u,\delta)$ is the smallest $u$ such that $\Pr(\mathrm{Bin}(n_u,u)\le n_u\widehat{q})\ge 1-\delta$.

\begin{remark}[Multiple SLAs and many-testing]
To maintain Certified~A's guarantees under updates, apply Bonferroni/Holm across all constraints and all candidate updates tested.
Disagreement gating reduces the \emph{data needed per update}, not the need for correct multiple-testing control.
\end{remark}

\section{Boundary audits for tail and stop-time SLAs}
Stop-time hiding and tail-latency guarantees are typically certified via rare-event indicators \cite{series:stoptime,series:scheduling}.
Naively estimating rare-event probabilities can require enormous samples.
However, anonymous DHT routers often produce a \emph{score} $s(X)$ (congestion/attack score) and implement policies as thresholds in $s$.
This induces a \emph{boundary} where violations concentrate.

A practical scheme:
\begin{itemize}
  \item bin contexts by score (or by proximity to policy thresholds);
  \item oversample and label more heavily near the boundary where tail events occur;
  \item combine bin-wise estimates into an overall tail certificate.
\end{itemize}

\begin{remark}[How to allocate boundary-audit labels across slices]
If contexts are stratified into slices $h$ with mass $p_h$ (from unlabeled telemetry) and within-slice standard deviation $s_h$ for the witness being certified, then under a fixed total label budget $n$ the variance of the stratified mean estimator is minimized by allocating labels as
$n_h \propto p_h s_h$ (Neyman/optimum allocation) \cite{neyman1934,cochran1977}.
In practice, use a pilot batch to estimate $s_h$ and round $n_h$ to integers.
\end{remark}



\paragraph{Off-policy and shadow-routing for audits.}
When the disagreement region is small, it can be efficient to log a small randomized exploration floor (shadow-routing) so that disagreement-band sampling is not biased.
The resulting audit is an off-policy evaluation problem; see Certified~A for background, references, and high-confidence off-policy certification pointers.

\section{Simulator evidence (kept in-\TeX)}
We include two small DHT-shaped tables to validate the workflow end-to-end.
As in Certified~A, these are not claimed as a full Octopus/Salsa/ShadowWalker simulator; they exist to make the certification and recertification protocol concrete.

\subsection{Disagreement gate: updates accepted with zero new labels}
We consider an OR-threshold anonymous router that chooses between two routing templates based on thresholds $(\theta_a,\theta_c)$ on an ``attack pressure'' score $a$ and ``congestion'' score $c$.
We estimate the disagreement rate $q$ from unlabeled contexts (here from Certified~A's simulator context distribution); Equation~\eqref{eq:disagree-area} gives a closed form under a uniform model as a sanity check.

Table~\ref{tab:gate-toy} shows three representative updates for a policy with privacy margin $\gamma=0.05$.
A small tweak passes the gate (no new labels); larger threshold changes fail the gate and trigger disagreement-band auditing.
The last column illustrates a simple label-budget heuristic: if a full-sample audit uses $n=12{,}000$ episodes, auditing only the disagreement region costs about $q n$ labels.

\begin{table}[t]
\centering
\begin{tabular}{@{}lrrr@{}}
\toprule
threshold update $(\theta_a,\theta_c)$ & disagreement $q$ & gate (margin $\gamma=0.05$) & labels if audit needed \\
\midrule
$(0.60,0.60)\to(0.61,0.59)$ & 0.0175 & pass & 0 \\
$(0.60,0.60)\to(0.65,0.65)$ & 0.0817 & fail & 981 \\
$(0.60,0.60)\to(0.40,0.40)$ & 0.3618 & fail & 4342 \\
\bottomrule
\end{tabular}
\caption{Disagreement gating for OR-threshold routers using an unlabeled context log (here drawn from Certified~A's simulator).
For the decision rule ``high iff $a>\theta_a$ or $c>\theta_c$,'' disagreement is the symmetric-difference area of two rectangles and can be computed exactly (Equation~\eqref{eq:disagree-area}).
If an update fails the gate, auditing only the disagreement region reduces labels from a full-sample audit to about $q$ times that amount.}
\label{tab:gate-toy}
\end{table}

\subsection{Boundary audits for rare tail constraints}
We illustrate boundary auditing for a rarer stop-time/latency constraint by slicing contexts by congestion quantiles.
Using the minimal simulator from Certified~A and a rarer tail threshold $t=1100$ms, tail violations concentrate in the high-congestion region.

\begin{table}[t]
\centering
\begin{tabular}{@{}lrrr@{}}
\toprule
congestion slice & slice mass & tail rate in slice & contribution to total tail \\
\midrule
bottom 90\% & 0.90 & $0.035778$ & $0.032200$ \\
next 7\% & 0.07 & $0.209714$ & $0.014680$ \\
next 2\% & 0.02 & $0.498000$ & $0.009960$ \\
top 1\% & 0.01 & $0.568000$ & $0.005680$ \\
\bottomrule
\end{tabular}
\caption{Boundary audit illustration from the minimal simulator (Certified~A), evaluated on one certified menu point, using a rarer tail threshold $t=1100$ms.
The unconditional tail rate is $0.0625$, while the conditional tail rate in the top 3\% congestion mass is $0.521$ (about $8.3\times$ larger), so oversampling near the boundary yields many more ``tail hits'' per label.}
\label{tab:boundary-sim}
\end{table}

\section{Exported object and earliest sufficient stop}
The stable output of this paper is narrower than a runtime receipt or a public release bundle.
A recertification packet needs only:
\begin{enumerate}[leftmargin=*]
  \item the ENF/TW/state identity verdict (same claim vs. new claim object),
  \item the disagreement estimate and its upper confidence bound,
  \item the zero-label gate verdict,
  \item if the gate fails, the disagreement-region conditional increment bound together with the labeling-propensity declaration, and
  \item if a tail or stop-time witness is near-tight, the published boundary slices and allocation rule used for the incremental audit.
\end{enumerate}
If a verifier only needs to decide whether a proposed update still lives under the old claim object and whether the old certificate still covers it (or a small incremental packet closes the gap), stopping at this packet is sufficient.
Plan-pinning / epoch-conformance receipts are intentionally deferred to Operational~B, and evaluator/public packaging are intentionally deferred to Eval~1 and ABOM/OINL~+~Release~A.\cite{series:operationalB,series:eval1,series:releaseA}

\paragraph{A minimal public recertification card.}
Table~\ref{tab:minimal-recert-card} is the smallest public packet later terse papers should need when the live question is simply \emph{did this update stay under the old certified claim surface, and if so by what gate?}

\begin{table}[t]
\centering
\small
\begin{tabular}{@{}p{0.21\linewidth}p{0.33\linewidth}p{0.38\linewidth}@{}}
\toprule
Field & What must be published & Why it is first-sufficient \\
\midrule
Claim-identity verdict & Same-claim / new-claim decision together with the compared ENF/TW/state identifiers & No reuse claim is meaningful until the paper says whether the witness surface actually stayed fixed. \\
Baseline anchor & Certified menu / witness version and the baseline certified margin $\gamma$ & The disagreement gate is only interpretable relative to one prior certificate. \\
Disagreement upper bound & The published $\overline q$ together with the unlabeled-context scope it covers & This is the zero-label quantity the gate consumes. \\
Gate verdict & Accept-from-old-cert / incremental-evidence-required verdict & This is the first reusable public answer for downstream notes. \\
Incremental packet (if needed) & Disagreement-band conditional increment bound and labeling-propensity declaration & If zero-label reuse fails, this is the minimal extra object needed to close the gap. \\
Boundary-audit rule (if needed) & Published slice rule for near-tight tail / stop-time witnesses & Rare-event SLAs need one explicit targeting rule rather than ad hoc prose. \\
Evidence pointer & Digest for the supporting unlabeled/labeled audit bundle and its request hook & Auditors need one exact packet to request, not the whole research archive. \\
\bottomrule
\end{tabular}
\caption{A minimal public recertification card. This is the non-decorative public answer later terse papers can quote directly when the live issue is change-control under a fixed certified interface rather than runtime conformance, evaluator notarization, or release lineage.}
\label{tab:minimal-recert-card}
\end{table}

This card is intentionally narrow: it fixes one claim-identity verdict, one inherited margin, one disagreement bound, and one explicit escalation route without silently re-owning runtime receipts, evaluator transcripts, or signed public release objects. The maintained worked cut grounds the downstream drill in 
\nolinkurl{example\_change\_control.json}, 
\nolinkurl{example\_compare\_report.json}, 
\nolinkurl{example\_drift\_cases.json}, and 
\nolinkurl{example\_release\_action\_matrix.json}, while explicitly leaving the separate 
\nolinkurl{bundle://worked-example/recertification} support-bundle route unmaterialized in this archive cut; this paper remains the first public owner of the packet's semantics.\cite{series:workedexample,series:operationalB,series:eval1,series:releaseA}

\paragraph{One compact recertification ladder.}
Table~\ref{tab:recert-ladder} states the smallest fail-closed cutover later terse papers should quote when the live issue is not the derivation of $\overline q$ or $\overline{\Delta}$, but simply \emph{which reuse state did this update land in?}

\begin{table}[t]
\centering
\small
\begin{tabular}{@{}p{0.19\linewidth}p{0.27\linewidth}p{0.48\linewidth}@{}}
\toprule
Case & Smallest public verdict & What must happen next \\
\midrule
ENF/TW/state ids changed & \textbf{new claim object required} & Old/new deltas are descriptive only; do not reuse the old certificate margin as a live gate. \\
Same ids and $\overline q \le \gamma$ & \textbf{reuse old certificate} & Quote the minimal recertification card only; zero new privacy labels are needed for this change-control step. \\
Same ids, $\overline q > \gamma$, and $\overline q\,\overline{\Delta} \le \gamma$ & \textbf{incremental packet closes gap} & Publish the disagreement-band increment bound, labeling-propensity declaration, and any boundary-slice rule for near-tight rare-event witnesses. \\
Same ids and $\overline q\,\overline{\Delta} > \gamma$ & \textbf{fail closed under old certificate} & Mint a new certified menu point or a new claim object rather than describing the update as covered by the old packet. \\
\bottomrule
\end{tabular}
\caption{A compact fail-closed recertification ladder. It exposes the smallest reusable cutover between zero-label reuse, incremental disagreement evidence, and full recertification.}
\label{tab:recert-ladder}
\end{table}

\paragraph{Same-card / refreshed-wrapper cutover.}
A later terse paper may keep citing the same recertification card only when the compared ENF/TW/state identifiers, inherited certified margin, unlabeled-scope declaration behind $\overline q$, gate verdict, and any already-declared incremental packet / boundary-slice rule are unchanged and the successor merely refreshes outer packaging (date stamp, local prose, ABOM/OINL wrapper, or downstream release pointer).
If any of those recertification rows changes---new ids, new unlabeled scope or $\overline q$, flipped gate verdict, changed disagreement-band increment, changed labeling-propensity declaration, or changed boundary-slice family---the old card becomes historical and a fresh recertification card is mandatory.
Runtime epoch receipts and signed public lineage remain downstream with Operational~B and Release~A rather than being smuggled back into this note as ``just wrapper churn.''\cite{series:operationalB,series:releaseA}

\paragraph{Tiny gate drill.}
Suppose the old certified menu point carried margin $\gamma=0.020$ and the unlabeled disagreement audit for a proposed update produced the upper bound $\overline q=0.013$ under the same ENF/TW/state identifiers.
Then the public gate verdict is immediately
\[
\overline q \le \gamma \quad\Rightarrow\quad 0.013 \le 0.020,
\]
so the update stays under the old certificate with \emph{zero new privacy labels}.
If a later update instead yields $\overline q=0.028$, the card must flip to \emph{incremental evidence required}; at that point later terse papers should still stop at the same recertification card unless they specifically need the disagreement-band estimate or the boundary-audit packet itself.

\paragraph{Tiny worked-bundle guard/classification drill.}
For the maintained worked bundle, the recertification-side guard is no longer hypothetical.
\texttt{example\_change\_control.json} fixes the same-interface guard at stable \texttt{tw\_id=sha256:c1ff81fc84049a2160a21487f61a4565e03c5263f1437649c4299d1354b35c18}, stable \texttt{state\_decl\_id=sha256:ddfbb55199a6ecdb225cbafc60a74c234ae3334fb40528fde0a2ec2ce6de3f20} for \texttt{stateadj.publicbucket.v1}, unchanged required exposures, and primary separation class \texttt{sep.nr1}, while \texttt{example\_compare\_report.json} records the concrete observed drift: \texttt{fallback\_contact\_surface} and \texttt{fallback\_tiered\_vector} changed because \texttt{obs\_model.p\_obs\_upper} decreased from \texttt{0.02} to \texttt{0.015}, so the maintained classification is \texttt{replay\_changed\_surfaces} and the required replay scope is exactly \texttt{eval-fallback-cct-v1} plus \texttt{eval-fallback-vector-v1}.
A later terse paper may therefore say only: \emph{the successor kept the certified interface guard and owed fresh replay of the changed fallback surfaces, not full recertification.}
That sentence is honest because Certified~C owns the guard/classification stop, the maintained source-object quartet carries the concrete guard fields, changed slices, and replay-versus-recertify cutover, and Release~A owns the downstream continuity/publication closure rather than having this note smuggle those release duties into the recertification card. The paper's own minimal public card is therefore paired with an explicit nonmaterialized-route boundary instead of standing in for a support bundle this cut does not ship.\cite{series:workedexample,series:releaseA}

\paragraph{One compact three-stop route.}
Suppose Certified~A already fixes the baseline menu card for witness/interface \texttt{cmenu-v1}, and suppose an optional Certified~B signature-manifest card separately fixes that a wider raw routing pool compresses from $M=12$ to $M_{\mathrm{eff}}=3$ on the declared weekday/weekend support.
This note owns only the third question: whether a deployed update still sits under that already-declared claim surface, for example because the same ENF/TW/state identifiers remain in force and the published disagreement upper bound satisfies $\overline q=0.013\le \gamma=0.020$.
If the live issue is only the baseline menu theorem, stop at Certified~A; if it is only why raw multiplicity was compressed, stop at Certified~B.
The stable certified-routing spine is therefore explicit: theorem/menu root in Certified~A, optional compression witness in Certified~B, and deployed change-control packet here. The maintained worked cut makes that third stop auditable through the four source objects above while still leaving 
\nolinkurl{bundle://worked-example/recertification} unmaterialized.\cite{series:menus,series:sigcomp,series:workedexample}

\paragraph{Tiny worked-bundle case ladder drill.}
The maintained drift catalog already exhibits the three fail-closed classes this paper wants later terse notes to quote, and it does so without decorative prose.
In \texttt{example\_drift\_cases.json}, \texttt{refresh-only-A} changes only the release-receipt checkpoint and remains \texttt{refresh\_only}; \texttt{replay-required-B} keeps the same \texttt{tw\_id} and \texttt{state\_decl\_id} but lowers \texttt{obs\_model.p\_obs\_upper} from \texttt{0.02} to \texttt{0.015}, so \texttt{example\_release\_action\_matrix.json} resolves it to rule \texttt{contact-or-budget-replay}; and \texttt{full-recertification-C} changes both \texttt{tw\_id} and \texttt{state\_decl\_id}, so the same action matrix resolves it to \texttt{interface-or-state-recertify}.
Later terse papers may therefore quote the cutover ladder in one line without rebuilding it: \emph{checkpoint-only drift refreshes the wrapper, fallback-surface drift replays the old certified interface, and interface/state drift reopens recertification.}
If a new case does not fit one of those maintained classes, reopen this paper rather than stretching the old ladder.

\paragraph{Tiny boundary-slice drill.}
Using Table~\ref{tab:boundary-sim}, the unconditional tail rate is $0.0625$ while the conditional tail rate in the top $3\%$ congestion mass is $0.521$.
A pilot batch of $200$ unconditional labels would therefore be expected to see about $200\cdot 0.0625 = 12.5$ tail hits, whereas $200$ labels concentrated in the declared boundary slice would be expected to see about $200\cdot 0.521 = 104.2$ tail hits.
That is the public reason the card's \emph{boundary-audit rule} may name one explicit slice family rather than vague ``extra tail checks'' prose: the slice rule is what turns a fail-open audit wish into a quotable, auditable targeting choice.
Later terse papers should still stop at the recertification card plus the ladder above unless they specifically need the full slice allocation schedule or the underlying labeled bundle.

\section{Handoff to runtime monitoring}

Between discrete recertifications, operators often want an \emph{anytime} alarm if SLAs begin to fail under drift. This paper does not claim ownership of those runtime-conformance receipts; it only specifies how a drift alarm should hand off into the recertification packet above. Operational~B owns the plan-pinning / epoch-receipt layer.\cite{series:operationalB}
Time-uniform confidence sequences for means provide such guarantees under minimal assumptions \cite{howard2021} and fit naturally into safe anytime-valid inference workflows \cite{ramdas2023savi}.
A simple deployment pattern:
\begin{itemize}
  \item maintain an always-valid upper confidence bound on $\mathbb{E}[W_{\pi_t}]$ from streaming telemetry;
  \item if the bound crosses $0$, trigger disagreement-gated recertification for a new policy.
\end{itemize}
This prevents optional stopping from invalidating risk monitoring.

\paragraph{Many updates over time.}
If operators test a stream of policy updates or monitor multiple menu points simultaneously, the relevant multiplicity grows over time.
Online family-wise error control (alpha-spending and its refinements) provides a principled way to maintain a global error budget across an unbounded sequence of tests \cite{tian2021onlinefwer}.
Equivalently, one can use \emph{e-values} / test supermartingales to obtain optional-stopping-safe alarms and combine evidence across time \cite{ramdas2025evals,vovk2021evals,hartog2025efwer}.


\section*{Artifacts}
Full scripts and raw sizing curves are in the research archive (available on request). This note keeps only the theory and minimal evidence.

\section{Takeaways}
\begin{itemize}[leftmargin=*]
\item \textbf{Same claim first:} disagreement-gated reuse is sound only after the ENF/TW/state surface is checked and declared unchanged (or explicitly coarsened).
\item \textbf{Zero-label updates are possible:} if disagreement stays below the old certificate margin, a policy update can be accepted without fresh privacy labels.
\item \textbf{Audit only where the update moves:} if the gate fails, label the disagreement region and concentrate rare-event effort near tail/stop-time boundaries.
\item \textbf{Stop at the recertification packet:} this paper is the certified-routing change-control companion only; runtime conformance receipts, evaluator notarization, and public release lineage belong to Operational~B, Eval~1, and Release~A rather than to this paper.
\item \textbf{Keep the exact public card honest:} this paper still owns the recertification card's semantics, and the maintained worked cut exposes the concrete guard/classification/case-ladder source objects without materializing a separate \nolinkurl{bundle://worked-example/recertification}; later terse papers should cite that boundary rather than improvising a fresh packet from memory.
\end{itemize}

\begin{thebibliography}{99}\setlength{\itemsep}{0pt}
\bibitem{neyman1934}
J.~Neyman.
On the Two Different Aspects of the Representative Method: The Method of Stratified Sampling and the Method of Purposive Selection.
\emph{Journal of the Royal Statistical Society}, 97(4):558--606, 1934.
DOI: 10.1111/j.2397-2335.1934.tb04184.x.

\bibitem{cochran1977}
W.~G. Cochran.
\emph{Sampling Techniques} (3rd ed.).
John Wiley \& Sons, 1977.


\bibitem{howard2021}
S.~R. Howard, A.~Ramdas, J.~McAuliffe, and J.~Sekhon.
Time-uniform, nonparametric, nonasymptotic confidence sequences.
\emph{Annals of Statistics}, 49(2):1051--1080, 2021.

\bibitem{clopper1934}
C.~J. Clopper and E.~S. Pearson.
The use of confidence or fiducial limits illustrated in the case of the binomial.
\emph{Biometrika}, 26(4):404--413, 1934. DOI: 10.1093/biomet/26.4.404.

\bibitem{ramdas2023savi}
A.~Ramdas, P.~Gr{\"u}nwald, V.~Vovk, and G.~Shafer.
Game-Theoretic Statistics and Safe Anytime-Valid Inference.
\emph{Statistical Science}, 38(4):576--597, 2023.
DOI: 10.1214/23-STS894. (Also arXiv:2210.01948.)

\bibitem{tian2021onlinefwer}
J.~Tian and A.~Ramdas.
Online control of the familywise error rate.
\emph{Statistical Methods in Medical Research}, 30(4):976--996, 2021. (Also arXiv:1910.04900.)

\bibitem{hartog2025efwer}
W.~Hartog and L.~Lei.
Family-wise Error Rate Control with E-values.
\emph{arXiv:2501.09015}, 2025. DOI: 10.48550/arXiv.2501.09015.

\bibitem{ramdas2025evals}
A.~Ramdas and R.~Wang.
Hypothesis testing with e-values.
\emph{Foundations and Trends in Statistics} (monograph); arXiv:2410.23614v6, 2025.
DOI: 10.48550/arXiv.2410.23614.

\bibitem{horvitz1952}
D.~G. Horvitz and D.~J. Thompson.
A generalization of sampling without replacement from a finite universe.
\emph{Journal of the American Statistical Association}, 47(260):663--685, 1952.
DOI: 10.1080/01621459.1952.10483446.

\bibitem{thomas2015hcope}
P.~S. Thomas, G.~Theocharous, and M.~Ghavamzadeh.
High-Confidence Off-Policy Evaluation.
\emph{AAAI Conference on Artificial Intelligence}, 2015, pp.~3000--3006.
\url{https://people.cs.umass.edu/~pthomas/papers/Thomas2015.pdf}.

\bibitem{laroche2019spibb}
R.~Laroche, P.~Trichelair, and R.~Tachet des~Combes.
Safe Policy Improvement with Baseline Bootstrapping.
\emph{ICML (PMLR~97)}, 2019.
\url{https://proceedings.mlr.press/v97/laroche19a/laroche19a.pdf}.

\bibitem{sridhar2024censorship}
S.~Sridhar, O.~Ascigil, N.~V.~Keizer, F.~Genon, S.~Pierre, Y.~Psaras, E.~Rivi\`ere, and M.~Kr\'ol.
Content Censorship in the InterPlanetary File System.
\emph{Network and Distributed System Security Symposium (NDSS)}, 2024.
DOI: 10.14722/ndss.2024.23153. (Also arXiv:2307.12212.)
\url{https://www.ndss-symposium.org/ndss-paper/content-censorship-in-the-interplanetary-file-system/}.

\bibitem{cholez2024sybilagain}
T.~Cholez and C.-L.~Ignat.
Sybil Attack Strikes Again: Denying Content Access in IPFS with a Single Computer.
\emph{Proceedings of the 19th International Conference on Availability, Reliability and Security (ARES)}, 2024.
DOI: 10.1145/3664476.3664482.
\url{https://dl.acm.org/doi/10.1145/3664476.3664482}.

\bibitem{netto2025activesybil}
V.~H.~M.~Netto, T.~Cholez, and C.-L.~Ignat.
Active Sybil Attack and Efficient Defense Strategy in IPFS DHT.
\emph{arXiv:2505.01139}, 2025. DOI: 10.48550/arXiv.2505.01139.

\bibitem{vovk2021evals}
V.~Vovk and R.~Wang.
E-values: Calibration, combination, and applications.
\emph{Annals of Statistics}, 49(3):1736--1754, 2021. DOI: 10.1214/20-AOS2020. (arXiv:1912.06116).

\bibitem{hanneke2014}
S.~Hanneke.
Theory of Disagreement-Based Active Learning.
\emph{Foundations and Trends in Machine Learning}, 7(2--3):131--309, 2014. DOI: 10.1561/2200000037.

\bibitem{series:menus}
Anonymous DHT Series.
\emph{Certified Menus for Anonymous DHT Lookups}.
(Certified~A in this series, draft.)

\bibitem{series:anytimespend}
Anonymous.
\newblock Anytime Statistical Spending for Menu Audits: Change-Control-Safe Confidence Budgets for Expanding Log Certificates.
\newblock Synthesis~26, The-One-True-Archive (2026).



\bibitem{series:artifactpolicy}
Anonymous.
\newblock Artifact and Disclosure Policy for the One-True-Archive (Source-Only Public Release).
\newblock Series synthesis note (Synthesis~6), 2026. (This archive.)

\bibitem{series:workedexample}
Anonymous.
\newblock Worked Example: Receipt Interlock for an Anonymous-DHT Claim Bundle.
\newblock Series synthesis note (Synthesis~17), 2026. (This archive.)

\bibitem{series:seriesspines}
Anonymous.
\newblock Series Spines for Anonymous-DHT Notes: Math-to-Review Handoffs.
\newblock Series synthesis note (Synthesis~55), 2026. (This archive.)

\bibitem{series:sigcomp}
Certified Series.
\emph{Routing-Signature Compression}.
(Series paper in this archive: \texttt{certified\_series/paperB\_routing\_signature\_compression/paper.tex}.)

\bibitem{series:scheduling}
Anonymous DHT Series.
\emph{Scheduling and Compiler Techniques for Metadata-Hiding Overlay Lookups}.
Wiki: \texttt{[[2026.01.22 - Mathematics: Scheduling and Compiler Techniques for Metadata-Hiding Overlay Lookups]]}.
(Series paper, already published.)

\bibitem{series:stoptime}
Anonymous DHT Series.
\emph{Stop-Time Padding Addendum (unique-only)}.
(Series addendum, already published.)

\bibitem{matter2025bgp}
J.~Matter and M.~Tran.
Network-level Censorship Attacks in the InterPlanetary File System.
\emph{arXiv:2509.06626}, 2025. DOI: 10.48550/arXiv.2509.06626.
\url{https://arxiv.org/abs/2509.06626}.

\bibitem{prunster2022disrupting}
B.~Pr\"unster, A.~Marsalek, and T.~Zefferer.
Total Eclipse of the Heart -- Disrupting the InterPlanetary File System.
\emph{31st USENIX Security Symposium (USENIX Security~'22)}, 2022.
\url{https://www.usenix.org/system/files/sec22-prunster.pdf}.

\bibitem{i2pnetdb}
I2P Project.
The Network Database (netDb) documentation (floodfill; XOR-metric distance; daily routing-key ``keyspace rotation'').
\url{https://beta.i2p.net/en/docs/overview/network-database/}.
(Updated Mar.~2025; accurate for 0.9.65; accessed Feb.~2026.)

\bibitem{i2pnetdbdiscussion}
I2P Project.
Network Database discussion (history; legacy Kademlia fallback disabled).
\url{https://geti2p.net/en/docs/discussions/netdb}.
(Accessed Feb.~2026.)

\bibitem{egger2013i2p}
C.~Egger, J.~Schlumberger, C.~Kruegel, and G.~Vigna.
Practical Attacks against the I2P Network.
\emph{Research in Attacks, Intrusions, and Defenses (RAID)} (LNCS~8145), 2013, pp.~432--451.
DOI: 10.1007/978-3-642-41284-4\_22.\;\url{https://link.springer.com/chapter/10.1007/978-3-642-41284-4_22}.

\bibitem{i2p2025timewilltell}
H.~Wang, Z.~Ling, X.~Xu, Y.~Pan, G.~Liu, J.~Luo, and X.~Fu.
\newblock Time will Tell: Large-scale De-anonymization of Hidden I2P Services via Live Behavior Alignment (Extended Version).
\newblock \emph{arXiv:2512.15510}, Dec.\ 2025.\;\url{https://arxiv.org/abs/2512.15510}.
\bibitem{i2p2026ndss}
H.~Wang, Z.~Ling, X.~Xu, Y.~Pan, G.~Liu, J.~Luo, and X.~Fu.
\newblock Time will Tell: Large-scale De-anonymization of Hidden I2P Services via Live Behavior Alignment.
\newblock In \emph{Network and Distributed System Security Symposium (NDSS)}, 2026.\;\href{https://www.ndss-symposium.org/ndss-paper/time-will-tell-large-scale-de-anonymization-of-hidden-i2p-services-via-live-behavior-alignment/}{NDSS 2026 paper page}.
\bibitem{wang2013kad}
P.~Wang, J.~Tyra, E.~Chan-Tin, T.~Malchow, D.~Foo~Kune, N.~Hopper, and Y.~Kim.
Attacking the Kad network---real world evaluation and high fidelity simulation using DVN.
\emph{Security and Communication Networks}, 6(12):1556--1575, 2013.
DOI: 10.1002/sec.172.


\bibitem{series:traceifaces}
Anonymous.
\newblock Trace Interfaces for Anonymous DHTs: Observation Tiers, Committee-Contact Traces, and Exposure Normal Forms.
\newblock Synthesis~3, The-One-True-Archive (2026).

\bibitem{series:receiptitems}
Anonymous.
\newblock Receipt Line Items for Anonymity Claims: A Minimal Tuple Schema for Published Budgets.
\newblock Synthesis~12, The-One-True-Archive (2026).

\bibitem{series:planstateequiv}
Anonymous.
\newblock Digest-Bound Replay Plans and State Declarations: Equivalence, Minimality, and Change Control.
\newblock Synthesis~18, The-One-True-Archive (2026).

\bibitem{series:contractobjects}
Anonymous.
\newblock Contract Objects for Anonymous-DHT Privacy Budgets and Claims.
\newblock Series synthesis note (Synthesis~0), 2026. (This archive.)

\bibitem{series:operationalB}
Anonymous.
\newblock Auditable Trace Privacy with Abort.
\newblock Operational series Paper~B, 2026. (This archive.)

\bibitem{series:eval1}
Anonymous.
\newblock Notarized Anonymity Certificates: Attempt Logs, Spending Discipline, and Verifier Rules.
\newblock Evaluation series Paper~1, 2026. (This archive.)

\bibitem{series:releaseA}
Anonymous.
\newblock Privacy as a Release Property for Anonymous DHTs.
\newblock Release and destination series Paper~A, 2026. (This archive.)

\bibitem{series:abomoinl}
Anonymous.
\newblock ABOM/OINL: Audit BOM and Outcome Identification Noise Label for Anonymous DHT Deployments.
\newblock Anonymity series Paper~C, 2026. (This archive.)

\end{thebibliography}

\end{document}
